% GENERATED FILE. Edit canonical lessons, then run npm run generate:books.
% canonical-source-hash: fnv1a64:f1fc430c1f3287e7
% canonical-lessons: ES-C674-pasear, ES-C674-paseo, ES-C674-callejear, ES-C674-recorrer, ES-C674-recorrido

\chapter{De paseo}
\label{ch:es-de-paseo}
\hlchaptermodality{674}

\begin{chapteropening}
\textbf{By the end of this chapter:} \emph{I can talk about walking round a town: strolling, going on foot and following a route.}

Complete the last lesson of chapter 674: I can talk about walking round a town: strolling, going on foot and following a route.
\end{chapteropening}

\section[pasear]{pasear — to go for a walk, to stroll}

\label{lesson:ES-C674-pasear}



\begin{quote}
Before the new one: say the Spanish for to go round, then the Spanish for to go along the edge of.
\end{quote}

\subsection*{You'll want to know: pasear}
\textbf{pasear} — \textquotedblleft{}to go for a walk, to stroll\textquotedblright{}.

\begin{quote}
Por la tarde paseamos por el puerto.
\end{quote}

\begin{quote}
\emph{In the evening we stroll round the harbour.}
\end{quote}

\emph{Pasear al perro} — to walk the dog.

\begin{cousinweb}
\emph{pasear} is built on \emph{paso} \textquotedblleft{}step\textquotedblright{} (EtymDB), which is inherited from Latin \emph{passus} \textquotedblleft{}a step\textquotedblright{}.
\end{cousinweb}

\begin{grammarlens}[title={What you've built}]
One more word for telling, explaining and describing at B1.
\end{grammarlens}

\subsection*{Your turn}
\begin{itemize}
\raggedright
  \item \emph{Say it:} \emph{pasear}
  \item \emph{Say it:} \emph{pasear}, once more
  \item \emph{Say it:} a sentence of your own with \emph{pasear}
  \item \emph{Recall:} say the Spanish for to go round, then the Spanish for to go along the edge of, then say \emph{pasear} again
\end{itemize}

\begin{tcolorbox}[breakable,colback=teal!4,colframe=teal!35!black,title={Before you move on}]
What does pasear mean? (\textquotedblleft{}To go for a walk, to stroll\textquotedblright{}.) Say it once more.
\end{tcolorbox}
\section[el paseo]{el paseo — the walk, the stroll; the promenade}

\label{lesson:ES-C674-paseo}



\begin{quote}
Before the new one: say the Spanish for to go along the edge of, then the Spanish for to go for a walk.
\end{quote}

\subsection*{You'll want to know: el paseo}
\textbf{el paseo} — \textquotedblleft{}the walk, the stroll; the promenade\textquotedblright{}.

\begin{quote}
Damos un paseo después de cenar.
\end{quote}

\begin{quote}
\emph{We go for a walk after dinner.}
\end{quote}

\emph{Dar un paseo} — to go for a walk. \emph{El paseo marítimo} — the seafront.

\begin{cousinweb}
\emph{paseo} is built on \emph{pasear} \textquotedblleft{}to stroll\textquotedblright{} (EtymDB), from \emph{paso} \textquotedblleft{}step\textquotedblright{}, Latin \emph{passus}.
\end{cousinweb}

\begin{grammarlens}[title={What you've built}]
One more word for telling, explaining and describing at B1.
\end{grammarlens}

\subsection*{Your turn}
\begin{itemize}
\raggedright
  \item \emph{Say it:} \emph{el paseo}
  \item \emph{Say it:} \emph{el paseo}, once more
  \item \emph{Say it:} a sentence of your own with \emph{el paseo}
  \item \emph{Recall:} say the Spanish for to go along the edge of, then the Spanish for to go for a walk, then say \emph{el paseo} again
\end{itemize}

\begin{tcolorbox}[breakable,colback=teal!4,colframe=teal!35!black,title={Before you move on}]
What does el paseo mean? (\textquotedblleft{}The walk, the stroll; the promenade\textquotedblright{}.) Say it once more.
\end{tcolorbox}
\section[callejear]{callejear — to wander the streets}

\label{lesson:ES-C674-callejear}



\begin{quote}
Before the new one: say the Spanish for to go for a walk, then the Spanish for the walk.
\end{quote}

\subsection*{You'll want to know: callejear}
\textbf{callejear} — \textquotedblleft{}to wander the streets\textquotedblright{}.

\begin{quote}
Me encanta callejear por el barrio viejo.
\end{quote}

\begin{quote}
\emph{I love wandering the streets of the old quarter.}
\end{quote}

Built on \emph{la calle}: walking with no plan.

\begin{cousinweb}
\emph{callejear} is built on \emph{calle} \textquotedblleft{}street\textquotedblright{}, which is inherited from Latin \emph{callis} \textquotedblleft{}a footpath\textquotedblright{}.
\end{cousinweb}

\begin{grammarlens}[title={What you've built}]
One more word for telling, explaining and describing at B1.
\end{grammarlens}

\subsection*{Your turn}
\begin{itemize}
\raggedright
  \item \emph{Say it:} \emph{callejear}
  \item \emph{Say it:} \emph{callejear}, once more
  \item \emph{Say it:} a sentence of your own with \emph{callejear}
  \item \emph{Recall:} say the Spanish for to go for a walk, then the Spanish for the walk, then say \emph{callejear} again
\end{itemize}

\begin{tcolorbox}[breakable,colback=teal!4,colframe=teal!35!black,title={Before you move on}]
What does callejear mean? (\textquotedblleft{}To wander the streets\textquotedblright{}.) Say it once more.
\end{tcolorbox}
\section[recorrer]{recorrer — to travel around, to go all over; to cover (a distance)}

\label{lesson:ES-C674-recorrer}



\begin{quote}
Before the new one: say the Spanish for the walk, then the Spanish for to wander the streets.
\end{quote}

\subsection*{You'll want to know: recorrer}
\textbf{recorrer} — \textquotedblleft{}to travel around, to go all over; to cover (a distance)\textquotedblright{}.

\begin{quote}
Queremos recorrer todo el país.
\end{quote}

\begin{quote}
\emph{We want to travel all over the country.}
\end{quote}

It takes a direct object: \emph{recorrer la ciudad}.

\begin{cousinweb}
\emph{recorrer} is inherited from Latin \emph{recurrere} \textquotedblleft{}to run back\textquotedblright{} (Lewis \& Short; EtymDB), built on \emph{currere} \textquotedblleft{}to run\textquotedblright{}.
\end{cousinweb}

\begin{grammarlens}[title={What you've built}]
One more word for telling, explaining and describing at B1.
\end{grammarlens}

\subsection*{Your turn}
\begin{itemize}
\raggedright
  \item \emph{Say it:} \emph{recorrer}
  \item \emph{Say it:} \emph{recorrer}, once more
  \item \emph{Say it:} a sentence of your own with \emph{recorrer}
  \item \emph{Recall:} say the Spanish for the walk, then the Spanish for to wander the streets, then say \emph{recorrer} again
\end{itemize}

\begin{tcolorbox}[breakable,colback=teal!4,colframe=teal!35!black,title={Before you move on}]
What does recorrer mean? (\textquotedblleft{}To travel around, to go all over; to cover (a distance)\textquotedblright{}.) Say it once more.
\end{tcolorbox}
\section[el recorrido]{el recorrido — the route, the tour}

\label{lesson:ES-C674-recorrido}



\begin{quote}
Before the new one: say the Spanish for to wander the streets, then the Spanish for to travel around.
\end{quote}

\subsection*{You'll want to know: el recorrido}
\textbf{el recorrido} — \textquotedblleft{}the route, the tour\textquotedblright{}.

\begin{quote}
El recorrido dura dos horas.
\end{quote}

\begin{quote}
\emph{The tour takes two hours.}
\end{quote}

\emph{Un recorrido guiado} — a guided tour.

\begin{cousinweb}
\emph{recorrido} is the past form of \emph{recorrer}, which is inherited from Latin \emph{recurrere} \textquotedblleft{}to run back\textquotedblright{}, built on \emph{currere} \textquotedblleft{}to run\textquotedblright{}.
\end{cousinweb}

\begin{grammarlens}[title={What you've built}]
One more word for telling, explaining and describing at B1.
\end{grammarlens}

\subsection*{Your turn}
\begin{itemize}
\raggedright
  \item \emph{Say it:} \emph{el recorrido}
  \item \emph{Say it:} \emph{el recorrido}, once more
  \item \emph{Say it:} a sentence of your own with \emph{el recorrido}
  \item \emph{Recall:} say the Spanish for to wander the streets, then the Spanish for to travel around, then say \emph{el recorrido} again
\end{itemize}

\begin{tcolorbox}[breakable,colback=teal!4,colframe=teal!35!black,title={Before you move on}]
What does el recorrido mean? (\textquotedblleft{}The route, the tour\textquotedblright{}.) Say it once more.
\end{tcolorbox}
